\subsection{Finite-Noise Performance}\label{sec:finite-noise}
We next evaluate the same encoders with the projection decoder of Section~\ref{sec:decoding}, followed by clipping to $[0,1]$. This comparison separates uniform disturbance robustness from average performance under a specified input distribution. For $p\in\{2,4\}$, write
$$
 f_p=t_p-m^p,\qquad
 \enc(s)=\sqrt P\,[as+i(bs^p+c)],\quad a,b>0.
$$
Let $N=\sqrt{N_0}Z$, where $Z$ is circular Gaussian with $\mathbb E|Z|^2=1$, independent of the sources. Thus each real noise component has variance $N_0/2$.

For independent sources with a common continuous density on $(-1,1)$, the target lies strictly inside $(0,1)$ almost surely. Proposition~\ref{prop:general-small-noise} therefore specializes to
\begin{equation}
 \lim_{N_0\downarrow0}\frac{\operatorname{MSE}(\enc,N_0)}{N_0}
 =\frac{1}{2K^2P}\left(\frac{1}{b^2}
 +\frac{p^2\mathbb E[m^{2p-2}]}{a^2}\right).
 \label{eq:small-noise-mse}
\end{equation}


The constrained theorem in Section~\ref{sec:pareto} interpolates between these objectives. The worst-case criterion uses the endpoint slope $M=p$, whereas \eqref{eq:small-noise-mse} uses $M_{\rm eff}=p\sqrt{\mathbb E[m^{2p-2}]}$. We construct a prior-aware comparison curve by inserting $M_{\rm eff}$ into the quadratic or quartic coefficient formulas. This choice minimizes the leading small-noise coefficient within the stated orthogonal curve family. We evaluate its moments analytically, without fitting simulation data, and use the same decoder for all designs. This comparison does not optimize finite-noise Bayesian risk.




We set $P=1$, $K\in\{4,16\}$, and draw independent uniform or endpoint-Beta sources. For each input distribution and transmitter count, we generate $200{,}000$ input vectors and independent standardized Gaussian disturbances. We reuse these draws across designs, targets, and the nine SNR values from $-10$ to $30$ dB in $5$-dB steps. Here SNR means the per-transmitter peak ratio $P/N_0$, not an aggregate received-power ratio. Paired standard errors account for the shared draws; intervals are pointwise, without a multiple-comparison adjustment.

\begin{table}[t]
\caption{MSE at $P/N_0=20$ dB, in units of $10^{-4}$. U: uniform; B: endpoint Beta. Robust: separation-optimal curve; tuned: prior-aware small-noise curve.}
\label{tab:finite-noise}
\centering\small
\begin{tabular}{@{}cccrrr@{}}\toprule
$p$ & Prior & $K$ & Baseline & Robust & Tuned\\\midrule
2 & U & 4  & 4.055 & 3.318 & 2.684\\
2 & B & 4  & 4.787 & 4.075 & 3.634\\
2 & U & 16 & 0.2115 & 0.1630 & 0.09922\\
2 & B & 16 & 0.2252 & 0.1772 & 0.1232\\
4 & U & 4  & 4.137 & 3.341 & 1.652\\
4 & B & 4  & 5.627 & 4.548 & 3.338\\
4 & U & 16 & 0.2450 & 0.1970 & 0.05521\\
4 & B & 16 & 0.2471 & 0.1987 & 0.06590\\\bottomrule
\end{tabular}
\end{table}

Across the 72 tested target/prior/$K$/SNR combinations, the separation-optimal curves reduce MSE by $6.12\%$--$22.92\%$ relative to the specified baselines. However, the prior-aware curves further reduce MSE by $5.69\%$--$72.01\%$ relative to the separation-optimal curves. Figure~\ref{fig:finite-noise} and Table~\ref{tab:finite-noise} display both outcomes. These average-performance gains sacrifice worst-case robustness: for the quartic gap with uniform sources and $K=16$, prior-aware tuning reduces $\Lambda_\sq/(K^2P)$ from $0.05853$ to $0.007273$. At low SNR, even the known prior-mean estimator can perform better: for uniform sources, $K=4$, variance, and $-10$ dB, its MSE is $0.02294$, compared with $0.1032$ for the separation-optimal curve and implemented decoder. These results quantify the operating behavior of specified designs; they establish neither universal MSE dominance nor Bayesian optimality. At $30$ dB, the endpoint-Beta, $K=4$ normalized errors still lie $3.49\%$--$7.19\%$ below their small-noise limits, beyond Monte Carlo uncertainty; boundary clipping remains relevant at this finite noise level. The code records every tested condition, paired uncertainty, analytic prior moments, seeds, and decoder cross-checks.
